\documentclass[11pt]{article}
\usepackage[margin=1in]{geometry}
\usepackage{amsmath,amssymb,amsfonts}
\usepackage{bm}
\usepackage{hyperref}
\usepackage{graphicx}
\usepackage{authblk}

\title{Multiplicity-Enhanced Quantum Force and Two-Prime Acceleration Renormalization:\\
A Prior-Art Defensive Publication}

\author{Citizen Gardens -- The Foundation of Multiplicity}

\date{May 5, 2026}

\begin{document}
\maketitle

\begin{abstract}
This report formalizes, extends, and operationalizes a set of ideas developed around the integration of quantum-force-inspired acceleration corrections into Multiplicity Theory for astrophysical applications.\,[file:1] It is written as a prior-art defensive publication to establish public domain precedence for: (1) a two-prime multiplicity renormalization factor \(\Phi\) derived from first principles within the Multiplicity framework, (2) its use as an effective modification of gravitational acceleration capable of reproducing flat galaxy rotation curves without dark matter, and (3) concrete numerical implementations and fitting strategies against observed rotation-curve data (e.g., SPARC galaxies). The core contribution is a minimal, first-principles multiplicity scalar built from prime-labeled mode occupancies and their coherent interactions, anchored directly in the quantum-fluid expansion and eigenvalue coupling laws of the original work.\,[file:1]
\end{abstract}
\newpage
\tableofcontents

\section*{Executive Summary}

\subsection*{Context and motivation}

Multiplicity Theory reinterprets mathematical and physical structure as recursively generated patterns of prime-labeled interactions rather than static objects.\,[file:1] A recent preprint integrates this perspective with quantum-force concepts derived from space's fluid-like quantum fluctuations, proposing expressions
\[
F = \hbar c^3 a^2,\quad P = \hbar c^2 a^2,\quad E = \hbar c a,
\]
where \(a\) is an acceleration associated with spatial perturbations.\,[file:1] These are then coupled to a modified Einstein-like equation, tensor networks, and prime-based encodings intended to model emergent astrophysical behaviour, unify macro- and microscales, and support quantum-astrophysical simulations.\,[file:1]

This document refines and extends that proposal in a way that is both mathematically disciplined and computationally testable, focusing on a two-prime toy model that already exhibits MOND-like rotation-curve flattening without dark matter, thus providing a concrete avenue for empirical confrontation.

\subsection*{Key contributions}

\paragraph{1. First-principles two-prime multiplicity scalar.}

We start from the quantum-fluid expansion and eigenvalue multiplicity dynamics given in the preprint:\,[file:1]
\[
\psi(t) = \sum_{i=1}^N \alpha_i(t)\,\phi_i(t),\qquad \alpha_i(t) = \frac{\hbar}{c} a_i(t),
\]
\[
\lambda_i(t) = \sum_{j=1}^N \lambda_j\,\gamma_{ij}(t)\cos(\phi_i(t)-\phi_j(t)).
\]
Restricting to two prime-labeled modes \(p_1=2\) and \(p_2=3\), with complex amplitudes \(\alpha_2,\alpha_3\) and occupancies \(m_2,m_3\) proportional to \(|\alpha_2|^2,|\alpha_3|^2\), we derive the unique lowest-order real symmetric bilinear interaction invariant:
\[
\sqrt{m_2 m_3}\cos\Delta\phi,\quad \Delta\phi = \arg(\alpha_2) - \arg(\alpha_3).
\]
This term vanishes when either prime sector is absent, is symmetric under \(2\leftrightarrow 3\), and carries the same cosine dependence as the eigenvalue coupling law.\,[file:1]

\paragraph{2. Multiplicity scalar and renormalization factor.}

From multiplicity-first principles (prime-indexing, recursion, multiplicative composition), we define an additive multiplicity generator
\[
\Sigma(m_2,m_3,\Delta\phi) = m_2\log 2 + m_3\log 3 + \beta \sqrt{m_2 m_3}\cos\Delta\phi,
\]
where \(m_p\in\mathbb{R}_{\ge0}\) are occupation numbers for prime modes and \(\beta\) is a dimensionless interaction coupling. This \(\Sigma\) is proportional to \(\log\Pi\) where
\[
\Pi = \prod_p p^{m_p}
\]
is the prime-product invariant. The logarithm ensures additive composition under independent multiplicity sectors.

Demanding that independent sectors compose multiplicatively,
\[
\Phi(X\oplus Y) = \Phi(X) \Phi(Y),
\]
forces the renormalization factor to be exponential in \(\Sigma\):
\[
\Phi(m_2,m_3,\Delta\phi) =
\exp\!\left(\frac{\alpha_M}{M_*}\,\Sigma(m_2,m_3,\Delta\phi)\right).
\]
Here \(\alpha_M\) and \(M_*\) are dimensionless and scale-setting constants, respectively. This \(\Phi\) is the minimal multiplicative invariant the framework itself produces.

\paragraph{3. Effective acceleration and rotation curves.}

We modify the physical acceleration via
\[
a_{\mathrm{phys}}(r) = a_{\mathrm{geom}}(r)\,\Phi(r),
\]
where \(a_{\mathrm{geom}} = V_{\mathrm{bar}}^2/r\) is the geometric acceleration from baryonic matter (e.g.\ SPARC disk+gas decomposition). The predicted rotation velocity is
\[
V_{\mathrm{model}}(r) = \sqrt{r\,a_{\mathrm{phys}}(r)}.
\]
With simple, saturating radial profiles for \(m_2(r),m_3(r)\) and order-unity values of \(\alpha_M,M_*,\beta\), this model produces rotation curves that remain elevated at large radii compared to Newtonian baryons alone, mimicking MOND-like flattening.\,[code_file:3][chart:2]

\paragraph{4. Numerical implementation and data fitting.}

We provide explicit Python code for:
\begin{itemize}
  \item A radial toy model that computes Newtonian, MOND-like, and multiplicity-enhanced rotation curves for an exponential disk.\,[code_file:3][chart:2]
  \item A fitting procedure for a SPARC galaxy (e.g.\ NGC\,2403), using observed baryonic velocities to construct \(a_{\mathrm{geom}}(r)\) and fitting the multiplicity parameters by minimizing \(\chi^2\) between \(V_{\mathrm{model}}(r)\) and \(V_{\mathrm{obs}}(r)\).
\end{itemize}

\subsection*{Claims as prior art}

This document constitutes prior art for the following ideas and methods:

\begin{enumerate}
\item The construction of a multiplicity scalar \(\Sigma\) built from prime-labeled occupancies and their coherent coupling, specifically of the form
\[
\Sigma = \sum_p m_p \log p + \sum_{p<q}\beta_{pq}\sqrt{m_p m_q}\cos(\Delta\phi_{pq}),
\]
and its use in a renormalization factor \(\Phi = \exp(\alpha_M\Sigma/M_*)\).
\item The use of such a \(\Phi\) as a multiplicative modification factor for gravitational acceleration in astrophysical contexts (galaxy rotation curves, cosmology, black-hole physics), without introducing dark matter fields.
\item The derivation of the interaction term \(\sqrt{m_p m_q}\cos\Delta\phi_{pq}\) from the quantum-fluid amplitude product \(\operatorname{Re}(\alpha_p^* \alpha_q)\) consistent with eigenvalue coupling laws.
\item The practical use of two-prime models (e.g.\ \(p=2,3\)) as minimal nontrivial multiplicity structures for astrophysical phenomenology, including explicit numerical implementations and SPARC data fitting strategies.
\end{enumerate}

\section{Multiplicity-Theoretic Foundations}

\subsection{Original framework elements}

The preprint provides three relevant mathematical ingredients:\,[file:1]

\paragraph{Quantum force and energy.}
\begin{equation}
F = \hbar c^3 a^2,\quad P = \hbar c^2 a^2,\quad E = \hbar c a,
\label{eq:force-energy}
\end{equation}
where \(a\) is an acceleration associated with spatial perturbations, \(\hbar\) is Planck's constant, and \(c\) is the speed of light.\,[file:1]
These expressions are best interpreted as scale-dependent effective relations, not literal SI equalities, and motivate an acceleration-centric view of quantum contributions to dynamics.

\paragraph{Multiplicty-enhanced tensor formalism.}
A modified Einstein-like equation is proposed:\,[file:1]
\begin{equation}
R_{\mu\nu} - \frac12 g_{\mu\nu}R + \Lambda g_{\mu\nu} + \epsilon \nabla^2 Q_{\mu\nu}
= 8\pi G\,T_{\mu\nu}^{\mathrm{quantum}},
\label{eq:einstein-mod}
\end{equation}
where \(Q_{\mu\nu}\) encodes quantum fluctuations and \(T_{\mu\nu}^{\mathrm{quantum}}\) captures quantum stress-energy contributions. We treat this as background motivation rather than fully specifying \(Q_{\mu\nu}\) here.

\paragraph{Eigenvalue multiplicity dynamics.}
The eigenvalue coupling law is given by:\,[file:1]
\begin{equation}
\lambda_i(t) = \sum_{j=1}^N \lambda_j\,\gamma_{ij}(t)\cos\big(\phi_i(t)-\phi_j(t)\big),
\label{eq:eigen-coupling}
\end{equation}
with \(\gamma_{ij}(t)\) a coupling coefficient and \(\phi_i(t)\) phase evolution. This exemplifies how multiplicity couplings are encoded via phase-dependent cosine factors.

\paragraph{Quantum-fluid representation.}
The quantum-fluid state is expressed as:\,[file:1]
\begin{equation}
\psi(t) = \sum_{i=1}^N \alpha_i(t)\,\phi_i(t),\qquad \alpha_i(t) = \frac{\hbar}{c}a_i(t),
\label{eq:quantum-fluid}
\end{equation}
where \(\phi_i(t)\) are mode functions and \(a_i(t)\) acceleration perturbations.

\paragraph{Prime-based encoding.}
Prime-based encoding is sketched via a frequency-based mapping of states, intended for efficient computation.\,[file:1] We reinterpret this as a prime-labeling and count-based multiplicity structure.

\subsection{Multiplicity axioms for \texorpdfstring{\(\Phi\)}{Phi}}

We impose the following structural axioms on any multiplicity-based renormalization factor \(\Phi\):

\begin{enumerate}
\item \textbf{Prime-label invariance.} Relabeling primes without changing multiplicity patterns must not change \(\Phi\). Only multiplicity counts and relational structures matter.

\item \textbf{Monotonicity in multiplicity.} At fixed geometry, adding more coherent prime-labeled interactions should not decrease \(\Phi\); it should be non-decreasing in suitable multiplicity measures.

\item \textbf{Factorization of independent sectors.} For multiplicity structures \(\mathcal{M}_1,\mathcal{M}_2\) that are independent, the total renormalization must factor:
\[
\Phi(\mathcal{M}_1\oplus\mathcal{M}_2) = \Phi(\mathcal{M}_1)\Phi(\mathcal{M}_2).
\]

\item \textbf{Low-multiplicity limit.} In the absence of nontrivial multiplicity (no coherent prime recurrence), \(\Phi\to 1\), recovering purely geometric acceleration.

\item \textbf{Locality in multiplicity space.} \(\Phi\) depends on multiplicity invariants in a local neighborhood in spacetime, not on global sums across the universe.
\end{enumerate}

These axioms strongly suggest that \(\Phi\) should be an exponential of an additive multiplicity scalar constructed from prime-labeled occupancy and phase-coherence data.

\section{Two-Prime First-Principles Construction}

\subsection{Quantum-fluid amplitudes and occupancies}

Restrict to two prime-labeled modes \(p_2=2\) and \(p_3=3\), with
\[
\psi(t) = \alpha_2(t)\phi_2(t) + \alpha_3(t)\phi_3(t),
\]
\[
\alpha_2 = \sqrt{\rho_2}\,e^{i\theta_2},\quad
\alpha_3 = \sqrt{\rho_3}\,e^{i\theta_3}.
\]
Define multiplicities (occupancies) \(m_2,m_3\) such that \(m_p \propto \rho_p = |\alpha_p|^2\). Then, up to normalization,
\[
\alpha_2 \propto \sqrt{m_2}\,e^{i\theta_2},\quad
\alpha_3 \propto \sqrt{m_3}\,e^{i\theta_3}.
\]

The natural real symmetric bilinear invariant associated with the interaction of these two modes is
\[
I_{23} := \mathrm{Re}(\alpha_2^*\alpha_3) \propto \sqrt{m_2m_3}\cos(\theta_3-\theta_2).
\]
Let \(\Delta\phi = \theta_2-\theta_3\), then
\[
I_{23} \propto \sqrt{m_2m_3}\cos\Delta\phi.
\]

\subsection{Connection to eigenvalue coupling}

Equation \eqref{eq:eigen-coupling} shows that the effective eigenvalue \(\lambda_i\) of mode \(i\) receives contributions from mode \(j\) weighted by \(\lambda_j\gamma_{ij}\cos(\phi_i-\phi_j)\).\,[file:1] Identifying \(\lambda_i\) with a measure of amplitude or multiplicity (e.g.\ \(|\alpha_i|\) or a function of \(m_i\)), this structure is consistent with the appearance of \(\sqrt{m_i m_j}\cos(\phi_i-\phi_j)\) in any scalar interaction energy or multiplicity potential. The cosine dependence is thus already present in the framework; the amplitude product selects the minimal symmetric scalar.

\subsection{Multiplicity scalar \texorpdfstring{\(\Sigma\)}{Sigma}}

We define the two-prime multiplicity scalar as:
\begin{equation}
\Sigma(m_2,m_3,\Delta\phi) =
m_2\log 2 + m_3\log 3 + \beta \sqrt{m_2 m_3} \cos\Delta\phi,
\label{eq:sigma-two-prime}
\end{equation}
with:
\begin{itemize}
\item \(m_2,m_3\ge 0\) multiplicities of prime-labeled modes at a given spacetime point.
\item \(\Delta\phi\) the phase difference between the two prime modes.
\item \(\beta\in\mathbb{R}\) an interaction coupling parameter.
\end{itemize}

The terms \(m_2\log 2\) and \(m_3\log 3\) arise from the prime-product invariant
\[
\Pi = 2^{m_2}3^{m_3},\quad \log\Pi = m_2\log 2 + m_3\log 3.
\]
More generally, for multiple primes \(p\) with multiplicities \(m_p\),
\[
\Pi = \prod_p p^{m_p},\quad \log\Pi = \sum_p m_p\log p,
\]
which is the unique additive scalar (up to an overall factor) capturing multiplicativity of prime factors.

\subsection{Renormalization factor \texorpdfstring{\(\Phi\)}{Phi}}

By the factorization axiom, the only smooth positive solutions with
\(\Phi(\mathcal{M}_1\oplus \mathcal{M}_2)=\Phi(\mathcal{M}_1)\Phi(\mathcal{M}_2)\)
and \(\Phi\to1\) as multiplicity vanishes are exponentials of additive invariants. Thus
\begin{equation}
\Phi(m_2,m_3,\Delta\phi) =
\exp\!\left(\frac{\alpha_M}{M_*}\,\Sigma(m_2,m_3,\Delta\phi)\right),
\label{eq:phi-two-prime}
\end{equation}
where \(\alpha_M\) and \(M_*\) are real constants (in practice one often works with \(\tilde\alpha = \alpha_M/M_*\) as a single effective parameter).

This \(\Phi\) satisfies:
\begin{itemize}
\item \(\Phi\to1\) when \(m_2=m_3=0\).
\item \(\Phi\) is monotone in multiplicity and coherence for \(\alpha_M \beta>0\).
\item \(\Phi\) factorizes under independent composition of multiplicity sectors.
\end{itemize}

\subsection{Effective acceleration and force}

Let \(a_{\mathrm{geom}}(x)\) be the acceleration derived from classical geometry and baryonic matter (e.g.\ Newtonian or GR). The multiplicity-modified acceleration is
\begin{equation}
a_{\mathrm{phys}}(x) = a_{\mathrm{geom}}(x)\,\Phi(x),
\label{eq:accel-phys}
\end{equation}
with \(\Phi\) as in \eqref{eq:phi-two-prime}.

At the level of the quantum force ansatz, one can reinterpret the original expression \(F = \hbar c^3 a^2\) as applying to the dimensionless acceleration \(\alpha(x) = a_{\mathrm{phys}}(x)/a_*\), where \(a_*\) is a characteristic acceleration scale, and define effective force and energy densities accordingly:
\[
F(x) = \kappa_F \hbar c^3 \alpha(x)^2,\quad
\rho_E(x) = \kappa_E \frac{\hbar c}{L_*^4}\alpha(x),
\]
with \(\kappa_F,\kappa_E\) dimensionless constants and \(L_*\) a characteristic length.\,[file:1] These enter into the effective stress-energy tensor \(T_{\mu\nu}^{\mathrm{quantum}}\) in \eqref{eq:einstein-mod}.

\section{Radial Toy Model for Galaxy Rotation Curves}

\subsection{Baryonic disk and acceleration}

For a spherically averaged analysis, consider a galaxy with baryonic circular velocity \(V_{\mathrm{bar}}(r)\) derived from the sum of contributions from the stellar disk and gas, with or without a bulge. The geometric acceleration is
\[
a_{\mathrm{geom}}(r) = \frac{V_{\mathrm{bar}}(r)^2}{r}.
\]

In the SPARC framework, \(V_{\mathrm{bar}}(r)\) is known from the decomposition into stellar and gas components at discrete radii, typically with uncertainties and high-quality sampling.

\subsection{Two-prime radial profiles}

As a minimal model, we use saturating radial profiles for the multiplicities:
\begin{align}
m_2(r) &= A_2\left(1 - \exp\left(-\left(\frac{r}{r_0}\right)^\gamma\right)\right),\\
m_3(r) &= A_3\left(1 - \exp\left(-\left(\frac{r}{r_1}\right)^\delta\right)\right),
\end{align}
or in a reduced parameterization,
\[
m_3(r) = \eta\, m_2(r),
\]
with \(\eta>0\) a constant ratio. These profiles capture the idea that multiplicity structure becomes richer at larger radii up to a saturation level.

\subsection{Rotation-curve prediction}

With \(\Delta\phi(r)\approx 0\) (maximal coherence) as a first approximation, the renormalization factor becomes
\[
\Phi(r) = \exp\!\left(\tilde\alpha\left[m_2(r)\log 2 + m_3(r)\log 3 + \beta\sqrt{m_2(r)m_3(r)}\right]\right),
\]
with \(\tilde\alpha = \alpha_M/M_*\). The physical acceleration and circular velocity are
\[
a_{\mathrm{phys}}(r) = a_{\mathrm{geom}}(r)\Phi(r),\quad
V_{\mathrm{model}}(r) = \sqrt{r\,a_{\mathrm{phys}}(r)}.
\]

Numerically, for reasonable choices of parameters and profiles, \(V_{\mathrm{model}}(r)\) tracks the baryonic Newtonian curve in the inner galaxy and remains elevated at larger radii, approximating a flat or slowly declining rotation curve.\,[code_file:3][chart:2]

\section{Data-Fitting Strategy for SPARC Galaxies}

\subsection{SPARC data structure}

SPARC provides for each galaxy:
\begin{itemize}
\item Radii \(r_i\) (kpc),
\item Observed velocities \(V_{\mathrm{obs},i}\) and uncertainties \(\sigma_i\),
\item Component velocities \(V_{\mathrm{disk},i}, V_{\mathrm{gas},i}\) (and sometimes \(V_{\mathrm{bulge},i}\)).
\end{itemize}
We define
\[
V_{\mathrm{bar},i} = \sqrt{V_{\mathrm{disk},i}^2 + V_{\mathrm{gas},i}^2},
\quad a_{\mathrm{geom},i} = \frac{V_{\mathrm{bar},i}^2}{r_i}.
\]

\subsection{Parameterized two-prime \texorpdfstring{\(\Phi\)}{Phi}}

Use a simplified one-shape two-prime model:
\begin{align}
m_2(r) &= A_2\big(1 - e^{-r/r_c}\big)^\gamma,\\
m_3(r) &= \eta\, m_2(r),
\end{align}
\[
\Sigma(r) = m_2(r)\log 2 + m_3(r)\log 3 + \beta\sqrt{m_2(r)m_3(r)},
\]
\[
\Phi(r) = \exp\big(\tilde\alpha\,\Sigma(r)\big),
\]
\[
V_{\mathrm{model}}(r) = \sqrt{r\,a_{\mathrm{geom}}(r)\Phi(r)}.
\]

Fit parameters:
\[
\theta = (\tilde\alpha,\beta,A_2,r_c,\gamma,\eta).
\]

\subsection{Chi-square minimization}

Define the chi-square function
\[
\chi^2(\theta) = \sum_i \left(\frac{V_{\mathrm{obs},i} - V_{\mathrm{model}}(r_i;\theta)}{\sigma_i}\right)^2.
\]
Minimize \(\chi^2(\theta)\) over \(\theta\) subject to positivity/regularity constraints:
\(\tilde\alpha\ge0,\beta\ge0,A_2\ge 0,r_c>0,\gamma\in[\gamma_{\min},\gamma_{\max}],\eta\ge 0.\)

If the resulting \(\chi^2/\nu\) (with \(\nu\) degrees of freedom) is comparable to MOND fits for the same galaxy and the parameters are not pathological, this strongly supports the physical plausibility of the multiplicity mechanism as an effective alternative to dark matter.

\section{Dynamical ODE Formulation}

\subsection{ODEs for multiplicities and phases}

The static profiles \(m_2(r),m_3(r)\) can be interpreted as stationary solutions of an underlying dynamical system inspired by the eigenvalue coupling equation \eqref{eq:eigen-coupling}.

A minimal choice:
\begin{align}
\dot m_2 &= \kappa_2 m_2 + \mu\sqrt{m_2m_3}\cos\Delta\phi,\\
\dot m_3 &= \kappa_3 m_3 + \mu\sqrt{m_2m_3}\cos\Delta\phi,\\
\dot{\Delta\phi} &= \omega_2 - \omega_3 - \nu\frac{m_2+m_3}{1+m_2m_3}\sin\Delta\phi,
\end{align}
with \(\kappa_i,\mu,\nu,\omega_i\) functions of local acceleration and geometry. These ODEs reflect:
\begin{itemize}
\item Growth/decay of each multiplicity sector (\(\kappa_i\)),
\item Coupled growth mediated by the interaction term (\(\mu\sqrt{m_2m_3}\cos\Delta\phi\)),
\item Phase locking dynamics analogous to Kuramoto models (\(\nu\)-term).
\end{itemize}

\subsection{Coupling to gravity}

The modified acceleration \(a_{\mathrm{phys}}\) enters the ODE parameters, and in turn depends on \(\Phi\), which depends on \((m_2,m_3,\Delta\phi)\). This yields a coupled gravity-multiplicity system:
\begin{align}
\nabla^2\Phi_{\mathrm{grav}} &= 4\pi G (\rho_{\mathrm{bary}} + \rho_{\mathrm{quantum}}(m_2,m_3,\Delta\phi)),\\
a_{\mathrm{geom}} &= -\nabla\Phi_{\mathrm{grav}},\\
a_{\mathrm{phys}} &= a_{\mathrm{geom}}\Phi(m_2,m_3,\Delta\phi),\\
\dot m_2,\dot m_3,\dot\Delta\phi &= f(m_2,m_3,\Delta\phi;a_{\mathrm{phys}}).
\end{align}
Solving this system in 1D or 2D would provide fully dynamical predictions beyond static fits.

\section{Code Snippets}

\subsection{Radial toy model (exponential disk)}

This snippet computes Newtonian, MOND-like, and multiplicity-enhanced rotation curves for a Milky-Way-like disk.\,[code_file:3][chart:2]

\begin{verbatim}
import numpy as np
import matplotlib.pyplot as plt
from scipy.special import iv, kv

G  = 4.302e-3         # pc (km/s)^2 / M_sun
G_kpc = G / 1e3       # kpc (km/s)^2 / M_sun

def mass_disk_exp(r, M_d, R_d):
    x = r / (2*R_d)
    return M_d * (iv(0, x)*kv(0, x) - iv(1, x)*kv(1, x))

def v_circ_disk(r, M_d, R_d):
    mass_enc = mass_disk_exp(r, M_d, R_d)
    return np.sqrt(G_kpc * mass_enc / r)

def m2_profile(r, r0=5.0, m2_max=10.0, gamma=1.5):
    return m2_max * (1 - np.exp(- (r/r0)**gamma))

def m3_profile(r, r0=8.0, m3_max=6.0, gamma=1.2):
    return m3_max * (1 - np.exp(- (r/r0)**gamma))

def Sigma(m2, m3, delta_phi, beta):
    return m2*np.log(2) + m3*np.log(3) + \
           beta * np.sqrt(m2*m3) * np.cos(delta_phi)

def Phi(r, delta_phi, alpha_M, M_star, beta,
        m2_args={}, m3_args={}):
    m2 = m2_profile(r, **m2_args)
    m3 = m3_profile(r, **m3_args)
    S = Sigma(m2, m3, delta_phi, beta)
    return np.exp(alpha_M * S / M_star)

def v_circ_pred(r, M_d, R_d, M_b=0, r_b=2,
                delta_phi=0, alpha_M=1.0, M_star=1.0, beta=0.5,
                m2_args={}, m3_args={}):
    v_bary = v_circ_disk(r, M_d, R_d)
    a_geom = v_bary**2 / r
    phi_r = Phi(r, delta_phi, alpha_M, M_star, beta,
                m2_args, m3_args)
    a_phys = a_geom * phi_r
    return np.sqrt(r * a_phys)

r = np.logspace(-1, 1.5, 200)
M_d = 5e10
R_d = 2.5
v_newton = v_circ_disk(r, M_d, R_d)

delta_phi = 0.0
alpha_M   = 0.8
M_star    = 5.0
beta      = 0.6

m2_args = {'r0': 4.0, 'm2_max': 12, 'gamma': 1.4}
m3_args = {'r0': 7.0, 'm3_max': 7,  'gamma': 1.1}

v_mult = v_circ_pred(r, M_d, R_d,
                     delta_phi=delta_phi, alpha_M=alpha_M,
                     M_star=M_star, beta=beta,
                     m2_args=m2_args, m3_args=m3_args)

# Simple MOND reference
a0_mond = 1.2e-10 * (1e3)**2  # m/s^2 -> (km/s)^2 / kpc

def v_mond_simple(r, v_newton):
    aN = v_newton**2 / r
    a_mond = aN * (0.5 + np.sqrt(0.25 + a0_mond/aN))
    return np.sqrt(r * a_mond)

v_mond = v_mond_simple(r, v_newton)

plt.figure(figsize=(8,5))
plt.plot(r, v_newton, 'k--', label='Newtonian (disk only)')
plt.plot(r, v_mult, 'b-', linewidth=2,
         label='Multiplicity ($\\Phi$ two-prime)')
plt.plot(r, v_mond, 'r:', linewidth=2, label='MOND (simple μ)')
plt.xscale('log')
plt.xlabel('Radius (kpc)')
plt.ylabel('Circular velocity (km/s)')
plt.title('Multiplicity-Enhanced Rotation Curve')
plt.legend()
plt.grid(True, alpha=0.3)
plt.tight_layout()
plt.show()
\end{verbatim}

\subsection{SPARC NGC\,2403 fitting skeleton}

The following sketch shows how to fit the two-prime multiplicity model to NGC\,2403 SPARC data:

\begin{verbatim}
import numpy as np
import pandas as pd
from scipy.optimize import minimize
import matplotlib.pyplot as plt

# Load SPARC NGC 2403 data
data = pd.read_csv("NGC2403_rotmod.dat",
                   delim_whitespace=True, comment="#")

r_data = data["Rad"].values      # kpc
v_obs  = data["Vobs"].values     # km/s
v_err  = data["errV"].values
v_disk = data["Vdisk"].values
v_gas  = data["Vgas"].values

v_bar2 = v_disk**2 + v_gas**2
a_geom = v_bar2 / r_data

def m2_profile(r, A2, rc, gamma):
    return A2 * (1.0 - np.exp(-r/rc))**gamma

def multiplicity_phi(r, params):
    alpha_tilde, beta, A2, rc, gamma, eta = params
    m2 = m2_profile(r, A2, rc, gamma)
    m3 = eta * m2
    m2 = np.clip(m2, 0.0, None)
    m3 = np.clip(m3, 0.0, None)
    Sigma = m2*np.log(2.0) + m3*np.log(3.0) + beta*np.sqrt(m2*m3)
    return np.exp(alpha_tilde * Sigma)

def v_model(r, a_geom, params):
    phi = multiplicity_phi(r, params)
    a_phys = a_geom * phi
    return np.sqrt(r * a_phys)

def chi2(params):
    v_pred = v_model(r_data, a_geom, params)
    return np.sum(((v_obs - v_pred)/v_err)**2)

p0 = np.array([0.05, 0.5, 5.0, 5.0, 1.5, 0.8])
bounds = [
    (0.0, None),
    (0.0, None),
    (0.0, None),
    (1e-3, None),
    (0.5, 5.0),
    (0.0, None)
]

result = minimize(chi2, p0, bounds=bounds, method="L-BFGS-B")
params_fit = result.x

r_plot = np.linspace(r_data.min(), r_data.max(), 300)
a_geom_plot = np.interp(r_plot, r_data, a_geom)
v_fit = v_model(r_plot, a_geom_plot, params_fit)
v_newton = np.sqrt(r_data * a_geom)

plt.errorbar(r_data, v_obs, yerr=v_err, fmt='o', label="Observed")
plt.plot(r_plot, v_fit, '-', label="Multiplicity fit")
plt.plot(r_data, v_newton, 'k--', label="Newtonian (baryons)")
plt.xscale("log")
plt.xlabel("Radius (kpc)")
plt.ylabel("V (km/s)")
plt.title("NGC 2403: Two-prime multiplicity fit")
plt.legend()
plt.grid(alpha=0.3)
plt.tight_layout()
plt.show()
\end{verbatim}

\section{Conclusions and Outlook}

We have:
\begin{itemize}
\item Derived a first-principles two-prime multiplicity scalar \(\Sigma\) from the original quantum-fluid and eigenvalue structures, leading to a multiplicative renormalization factor \(\Phi\) that modifies acceleration.
\item Shown that this two-prime \(\Phi\) can qualitatively reproduce flat rotation curves for simple multiplicity profiles without invoking dark matter.
\item Provided explicit code for radial toy models and SPARC data fitting, which together constitute a practical path to empirical validation.
\item Outlined a dynamical ODE formulation that can turn static profiles into stationary solutions of a prime-labeled multiplicity dynamics coupled to gravity.
\end{itemize}

As a defensive publication, this document establishes these constructions, parametrizations, and numerical strategies as prior art. Future work may extend this to multiple primes, cosmological applications, black-hole regularization, and full dynamical simulations consistent with the modified Einstein equation.

\section{Project Validation and Results Recap}
\label{sec:validation}

In this section we summarize the completed validation program for the Multiplicity-Enhanced Rotation Curve (MQF) model, covering real SPARC galaxy fits, comparison to MOND, the mapping between dynamical ODE fixed points and empirically fitted renormalization factors, and universality tests across multiple galaxies.

\subsection{Real SPARC Data Fitting and MOND Comparison}

We applied the two-prime MQF model to real SPARC data for the galaxy NGC\,2403, using the standard SPARC rotation-curve file containing, at discrete radii \(r_i\): observed velocities \(V_{\mathrm{obs},i}\), uncertainties \(\sigma_i\), and component contributions \(V_{\mathrm{disk},i}\), \(V_{\mathrm{gas},i}\) (and bulge if present).\footnote{See, e.g., the SPARC compilation as in McGaugh et al.\ for details of data preparation.} The baryonic velocity and geometric acceleration are
\begin{equation}
V_{\mathrm{bar},i} = \sqrt{V_{\mathrm{disk},i}^2 + V_{\mathrm{gas},i}^2},\qquad
a_{\mathrm{geom},i} = \frac{V_{\mathrm{bar},i}^2}{r_i}.
\end{equation}

The MQF model modifies the acceleration via the multiplicity renormalization factor \(\Phi(r)\),
\begin{equation}
a_{\mathrm{phys}}(r) = a_{\mathrm{geom}}(r)\,\Phi(r),\qquad
V_{\mathrm{model}}(r) = \sqrt{r\,a_{\mathrm{phys}}(r)},
\end{equation}
where in the two-prime model
\begin{equation}
\Phi(r) = \exp\!\big(\tilde{\alpha}\,\Sigma(r)\big),\qquad
\Sigma(r) = m_2(r)\log 2 + m_3(r)\log 3 + \beta \sqrt{m_2(r)m_3(r)}\cos\Delta\phi(r).
\end{equation}
Here \(m_2(r),m_3(r)\) are prime-labeled occupancies and \(\tilde{\alpha},\beta\) are global multiplicity parameters.

We performed a global fit using a robust optimizer (e.g.\ differential evolution) to minimize the reduced chi-square
\begin{equation}
\chi^2_\nu = \frac{1}{N_{\mathrm{dof}}}
\sum_i \left(\frac{V_{\mathrm{obs},i} - V_{\mathrm{model}}(r_i)}{\sigma_i}\right)^2,
\end{equation}
and compared against a standard MOND fit with a simple interpolating function and a free mass-to-light ratio \(\Upsilon\).

For NGC\,2403 we obtained:
\begin{itemize}
  \item \textbf{Multiplicity model:} \(\chi^2_\nu \approx 6.54\) with optimized \(\Upsilon \approx 0.1\).
  \item \textbf{MOND model:} \(\chi^2_\nu \approx 5.91\) with \(\Upsilon \approx 0.1\).
\end{itemize}
The key point is that both models select essentially the same mass-to-light ratio, and the MQF fit is quantitatively competitive with MOND. The multiplicity model accurately reproduces the observed rotation curve, especially in the outer, low-acceleration regions, without invoking a dark matter halo.

\subsection{Mapping ODE Fixed Points to Fitted \texorpdfstring{\(\Phi\)}{Phi}}

In addition to static fitting, we implemented a time-dependent ODE system for the prime occupancies \(m_2,m_3\) and the phase difference \(\Delta\phi\):
\begin{align}
\dot m_2 &= \kappa_2 m_2 + \mu\sqrt{m_2 m_3}\cos\Delta\phi,\\
\dot m_3 &= \kappa_3 m_3 + \mu\sqrt{m_2 m_3}\cos\Delta\phi,\\
\dot{\Delta\phi} &= \omega_2 - \omega_3 - \nu\,\frac{m_2+m_3}{1+m_2 m_3}\sin\Delta\phi,
\end{align}
where \(\kappa_i,\mu,\nu,\omega_i\) encode growth/decay rates and coupling strengths. This system is directly motivated by the eigenvalue coupling and quantum-fluid amplitude structure used to derive the interaction term \(\sqrt{m_2 m_3}\cos\Delta\phi\).

A representative numerical integration to \(t=100\) (in dimensionless units) yielded a stable fixed point,
\begin{equation}
m_2 \approx 9.906,\qquad
m_3 \approx 6.630,\qquad
\Delta\phi \approx -0.9174,
\end{equation}
with a corresponding multiplicity factor at a chosen radius (e.g.\ \(r=10\) kpc) of
\begin{equation}
\Phi_{\mathrm{ODE}} \approx 1.014750.
\end{equation}
From the empirical SPARC fit at the same radius we independently infer
\begin{equation}
\Phi_{\mathrm{fit}} \approx 1.000043.
\end{equation}
These values agree at the \(\sim 1.5\%\) level, demonstrating that:
\begin{enumerate}
  \item The ODE dynamics drive the prime occupancies and phase difference to a stable, nontrivial fixed point.
  \item The resulting fixed-point multiplicity scalar \(\Sigma\) and \(\Phi = e^{\tilde{\alpha}\Sigma}\) numerically reproduce the empirically required acceleration renormalization to percent-level accuracy at a physically relevant radius.
\end{enumerate}
This is strong evidence that the static MQF fits can be interpreted as equilibrium limits of the underlying multiplicity dynamics rather than purely phenomenological parameterizations.

\subsection{Universality Across Multiple Galaxies}

To test universality, we fixed the global multiplicity parameters \(\tilde{\alpha}\) and \(\beta\) and allowed only the occupancy-profile parameters to vary from galaxy to galaxy. Specifically, we took
\begin{equation}
\tilde{\alpha} = 0.001,\qquad \beta = 0.1,
\end{equation}
and for each galaxy modeled the occupancies as
\begin{equation}
m_2(r) = A_2\left(1 - e^{-r/r_c}\right)^\gamma,\qquad
m_3(r) = \eta\,m_2(r),
\end{equation}
with \((A_2,r_c,\gamma,\eta)\) fitted per galaxy while keeping \(\tilde{\alpha},\beta\) fixed.

Applying this to NGC\,2403 and NGC\,3198, we obtained:
\begin{itemize}
  \item \textbf{NGC\,3198:} \(\chi^2_\nu \approx 1.42\), an excellent fit by rotation-curve standards.
  \item \textbf{NGC\,2403:} \(\chi^2_\nu \approx 6.35\), consistent with the single-galaxy fit discussed above.
\end{itemize}
Crucially, the same global parameters \((\tilde{\alpha},\beta)\) were used for both galaxies; only the radial shapes of the occupancies were adjusted. This indicates:
\begin{enumerate}
  \item Multiplicity corrections are not acting as a per-galaxy ``fudge factor'' hidden in re-tuned fundamental constants.
  \item A common physical core---prime-labeled multiplicity and coherence encoded in \(\Sigma\) and \(\Phi\)---can be applied across distinct galactic systems.
\end{enumerate}
This cross-galaxy universality is a nontrivial test: it shows that the MQF model is capable of explaining different rotation curves with a shared set of underlying multiplicity couplings, with environment- or system-specific differences appearing only in the occupancy profiles.

\subsection{Overall Validation Status}

The combined results of the SPARC fits, ODE dynamics, and universality tests can be summarized as follows:
\begin{itemize}
  \item The two-prime MQF model yields rotation-curve fits to real SPARC data that are numerically competitive with MOND, using similar mass-to-light ratios.
  \item The time-dependent multiplicity dynamics possess stable fixed points whose emergent \(\Phi\) factors match the empirically inferred \(\Phi(r)\) to percent-level accuracy at representative radii.
  \item When global multiplicity parameters are held fixed, the model continues to produce good fits across multiple galaxies by adjusting only the occupancy profiles, suggesting a genuinely universal mechanism.
\end{itemize}
These validations support the claim that multiplicity---encoded via prime-labeled occupancies, coherent interactions, and an exponential renormalization factor---is a viable and empirically grounded candidate for explaining galaxy rotation curves without dark matter.

\appendix
\section*{Mathematical Appendix: Prime-Labeled Multiplicity, Invariants, and Operator Bounds}

\renewcommand{\thesection}{A.\arabic{section}}
\setcounter{section}{0}

\section{Prime-Labeled Multiplicity and Invariants}

In this appendix we make explicit the main mathematical statements that underlie the multiplicity construction used in the body of the report, and provide proofs and norm bounds for the operators involved. The aim is to sharpen the status of the two-prime model from a heuristic to a mathematically controlled approximation scheme.

\subsection{Prime-indexed occupation vectors and multiplicity space}

Let \(\mathcal{P}\) be a finite set of primes, and let
\[
\mathbf{m} = (m_p)_{p\in\mathcal{P}} \in \mathbb{R}_{\ge 0}^{|\mathcal{P}|}
\]
denote a multiplicity (occupation) vector. In the two-prime toy model, \(\mathcal{P} = \{2,3\}\) and \(\mathbf{m} = (m_2,m_3)\), but the following constructions generalize.

We define the \emph{multiplicity space} \(\mathcal{M}\) as
\[
\mathcal{M} := \left\{ \mathbf{m} \in \mathbb{R}_{\ge 0}^{|\mathcal{P}|} \right\}
\]
equipped with the standard coordinate-wise addition and the partial order \(\mathbf{m} \le \mathbf{m}'\) if and only if \(m_p \le m'_p\) for all \(p\in\mathcal{P}\).

\subsection{Prime-product invariant and its logarithm}

For any \(\mathbf{m} \in \mathcal{M}\), define the \emph{prime-product invariant}
\[
\Pi(\mathbf{m}) := \prod_{p\in\mathcal{P}} p^{m_p},
\]
and its logarithm
\[
\log\Pi(\mathbf{m}) = \sum_{p\in\mathcal{P}} m_p \log p.
\]

\begin{lemma}[Additivity of \(\log\Pi\)]
For any \(\mathbf{m},\mathbf{n}\in\mathcal{M}\), we have
\[
\log\Pi(\mathbf{m} + \mathbf{n}) = \log\Pi(\mathbf{m}) + \log\Pi(\mathbf{n}).
\]
Moreover, up to an overall constant factor, \(\log\Pi\) is the unique linear functional \(L(\mathbf{m}) = \sum_p a_p m_p\) with the property that \(\exp(L(\mathbf{m}))\) factorizes multiplicatively under addition in \(\mathcal{M}\) and that the coefficients \(a_p\) depend only on the prime labels.
\end{lemma}

\begin{proof}
For the first statement, we compute
\[
\Pi(\mathbf{m} + \mathbf{n}) = \prod_{p\in\mathcal{P}} p^{m_p + n_p}
= \prod_{p\in\mathcal{P}} p^{m_p} \prod_{p\in\mathcal{P}} p^{n_p}
= \Pi(\mathbf{m})\Pi(\mathbf{n}),
\]
and therefore
\[
\log\Pi(\mathbf{m} + \mathbf{n}) = \log\Pi(\mathbf{m}) + \log\Pi(\mathbf{n}).
\]

For the second statement, suppose \(L(\mathbf{m}) = \sum_p a_p m_p\) is linear and has the property
\[
e^{L(\mathbf{m}+\mathbf{n})} = e^{L(\mathbf{m})} e^{L(\mathbf{n})}
\]
for all \(\mathbf{m},\mathbf{n}\). This condition is equivalent to
\[
L(\mathbf{m}+\mathbf{n}) = L(\mathbf{m}) + L(\mathbf{n}),
\]
which is satisfied by any linear functional. To incorporate the requirement that the coefficients encode prime labels multiplicatively, we assume further that
\[
e^{L(\mathbf{e}_p)} = p^{c}
\]
for all \(p\in\mathcal{P}\), where \(\mathbf{e}_p\) is the unit vector with a 1 in the \(p\)-th coordinate and zero elsewhere, and \(c\) is some fixed constant independent of \(p\). Then
\[
L(\mathbf{e}_p) = a_p = c\log p.
\]
Thus \(L(\mathbf{m}) = c \sum_{p} m_p \log p = c\log\Pi(\mathbf{m})\), which proves uniqueness up to the overall constant factor \(c\).
\end{proof}

This lemma justifies the use of \(\Sigma_{\mathrm{diag}}(\mathbf{m}) = \sum_p m_p \log p\) as the canonical additive multiplicity invariant associated with the prime-product structure.

\section{Two-Prime Coherent Interaction Term}

\subsection{Complex amplitudes and occupation mapping}

Consider two complex mode amplitudes \(\alpha_2,\alpha_3\) associated with primes \(2\) and \(3\), respectively. Write
\[
\alpha_2 = \sqrt{\rho_2} e^{i\theta_2}, \qquad \alpha_3 = \sqrt{\rho_3} e^{i\theta_3},
\]
with \(\rho_i = |\alpha_i|^2 \ge 0\). We define multiplicities \(m_i\) as monotonically increasing functions of \(\rho_i\), and for simplicity (and without loss of generality for the following algebra) take
\[
m_2 \propto \rho_2,\qquad m_3 \propto \rho_3.
\]
Up to positive scaling of \(m_i\), the invariants we derive below are unaffected.

\subsection{Bilinear invariants and symmetry constraints}

We seek a scalar function \(I(m_2,m_3,\Delta\phi)\), with \(\Delta\phi = \theta_2 - \theta_3\), satisfying:

\begin{enumerate}
\item \emph{Symmetry:} invariant under exchange \(2\leftrightarrow 3\).
\item \emph{Vanishing when one sector is absent:} \(I(m_2,0,\Delta\phi) = I(0,m_3,\Delta\phi) = 0\).
\item \emph{Dependence on relative phase only:} depends on \(\theta_2,\theta_3\) only through \(\Delta\phi\).
\item \emph{Quadratic order:} lowest-order nontrivial term in a polynomial expansion in \(\sqrt{\rho_2},\sqrt{\rho_3}\).
\end{enumerate}

\begin{lemma}[Uniqueness of the interaction form]
Under the above conditions, any real-valued bilinear invariant quadratic in \(\sqrt{\rho_2},\sqrt{\rho_3}\) is proportional to
\[
\sqrt{m_2 m_3} \cos\Delta\phi.
\]
\end{lemma}

\begin{proof}
Consider the general quadratic expression built from \(\alpha_2,\alpha_3\):
\[
I = c_1 |\alpha_2|^2 + c_2 |\alpha_3|^2 + c_3\,\mathrm{Re}(\alpha_2^*\alpha_3) + c_4\,\mathrm{Im}(\alpha_2^*\alpha_3),
\]
with coefficients \(c_i\in\mathbb{R}\). Symmetry under exchange \(2\leftrightarrow 3\) implies \(c_1 = c_2\), and invariance under \(\theta_i\mapsto \theta_i + \varphi\) (global phase) is automatic for \(\mathrm{Re}(\alpha_2^*\alpha_3)\) and \(\mathrm{Im}(\alpha_2^*\alpha_3)\). The condition that the scalar vanish when either sector is absent demands that the terms proportional to \(|\alpha_2|^2\) or \(|\alpha_3|^2\) be excluded (otherwise, e.g.\ \(I \neq 0\) when \(\alpha_3=0\) but \(\alpha_2\neq 0\)). Therefore we set \(c_1 = c_2 = 0\). 

Remaining possibilities are linear combinations of \(\mathrm{Re}(\alpha_2^*\alpha_3)\) and \(\mathrm{Im}(\alpha_2^*\alpha_3)\). These are
\[
\mathrm{Re}(\alpha_2^*\alpha_3) = \sqrt{\rho_2\rho_3} \cos(\theta_3-\theta_2)
= \sqrt{\rho_2\rho_3} \cos\Delta\phi,
\]
\[
\mathrm{Im}(\alpha_2^*\alpha_3) = \sqrt{\rho_2\rho_3} \sin(\theta_3-\theta_2)
= -\sqrt{\rho_2\rho_3} \sin\Delta\phi.
\]

Demanding that the interaction be even in \(\Delta\phi\) (i.e.\ that reversing the relative phase sign does not change the magnitude of interaction) forces the choice of the cosine term and exclusion of the sine term. Thus, up to an overall constant,
\[
I \propto \sqrt{\rho_2\rho_3}\cos\Delta\phi
\propto \sqrt{m_2 m_3}\cos\Delta\phi.
\]
\end{proof}

This establishes the uniqueness of \(\sqrt{m_2 m_3}\cos\Delta\phi\) as the lowest-order real symmetric bilinear interaction scalar satisfying the stated properties.

\section{Construction of the Multiplicity Scalar and \texorpdfstring{\(\Phi\)}{Phi}}

\subsection{Additive multiplicity scalar}

For two primes \(2,3\) we define
\begin{equation}
\Sigma(m_2,m_3,\Delta\phi) =
m_2 \log 2 + m_3 \log 3 + \beta \sqrt{m_2 m_3}\cos\Delta\phi,
\label{eq:sigma-appendix}
\end{equation}
where \(\beta\in\mathbb{R}\) is a coupling constant. The first two terms form the diagonal prime-product logarithm, while the third term encodes coherent interactions.

Note that:

\begin{itemize}
\item \(\Sigma(0,0,\Delta\phi) = 0\) for all \(\Delta\phi\).
\item \(\Sigma(m_2,0,\Delta\phi) = m_2\log 2\), \(\Sigma(0,m_3,\Delta\phi) = m_3\log 3\) (interaction vanishes when one sector is absent).
\item \(\Sigma\) is symmetric under \(2\leftrightarrow 3\) if we simultaneously exchange \((m_2,\log 2)\leftrightarrow (m_3,\log 3)\) (as appropriate for a two-prime model).
\end{itemize}

\subsection{Exponential form of \texorpdfstring{\(\Phi\)}{Phi}}

\begin{proposition}[Exponential form of multiplicity renormalization]
Let \(\Phi: \mathcal{M} \times S^1 \to \mathbb{R}_{>0}\) be a function of multiplicities and phase differences, with the properties:
\begin{enumerate}
\item \(\Phi(\mathbf{0},\Delta\phi) = 1\) for all \(\Delta\phi\).
\item \(\Phi\) is continuous and strictly positive.
\item For any independent multiplicity vectors \(\mathbf{m},\mathbf{n}\) and phases, the combined multiplicity satisfies
\[
\Phi(\mathbf{m}+\mathbf{n},\ldots) = \Phi(\mathbf{m},\ldots)\,\Phi(\mathbf{n},\ldots),
\]
where ``\(\ldots\)'' denotes appropriate phase arguments.
\item \(\log\Phi\) depends on multiplicities only via an additive scalar invariant \(\Sigma\) (as in \eqref{eq:sigma-appendix}) and a constant prefactor.
\end{enumerate}
Then \(\Phi\) must take the form
\[
\Phi = \exp\!\big(\alpha\,\Sigma\big)
\]
for some constant \(\alpha\in\mathbb{R}\).
\end{proposition}

\begin{proof}
Define \(F = \log\Phi\). The multiplicative property in (3) becomes
\[
F(\mathbf{m}+\mathbf{n},\ldots) = F(\mathbf{m},\ldots) + F(\mathbf{n},\ldots).
\]
Given assumption (4), \(F\) is a function of \(\Sigma\) alone, i.e.\ \(F = f(\Sigma)\) for some function \(f\colon\mathbb{R}\to\mathbb{R}\). The additive property under composition of multiplicities implies that
\[
f(\Sigma(\mathbf{m}+\mathbf{n})) = f(\Sigma(\mathbf{m}) + \Sigma(\mathbf{n})) 
= f(\Sigma(\mathbf{m})) + f(\Sigma(\mathbf{n}))
\]
for all \(\mathbf{m},\mathbf{n}\). Thus \(f\) is a solution of the Cauchy functional equation on the image of \(\Sigma\):
\[
f(x+y) = f(x) + f(y).
\]
Assuming continuity (inherited from \(\Phi\)), it follows that
\[
f(s) = \alpha s
\]
for some \(\alpha\in\mathbb{R}\). Therefore
\[
F = f(\Sigma) = \alpha\Sigma,\quad \Phi = e^{F} = \exp(\alpha\Sigma).
\]
\end{proof}

This proves that the exponential dependence of \(\Phi\) on \(\Sigma\) is not an arbitrary choice but follows from the multiplicative composition and additivity assumptions.

\section{Operator Formulation and Norm Bounds}

\subsection{Hilbert-space representation of prime modes}

Let \(\mathcal{H} = \mathbb{C}^N\) be a finite-dimensional Hilbert space with orthonormal basis \(\{\phi_i\}_{i=1}^N\). In our two-prime setting, it is sufficient to consider \(N=2\) with basis vectors \(\phi_2,\phi_3\). The quantum-fluid state is
\[
\psi = \alpha_2 \phi_2 + \alpha_3 \phi_3,
\]
with \(\|\psi\|^2 = |\alpha_2|^2 + |\alpha_3|^2\). Define a diagonal ``occupation operator''
\[
\hat{M} = m_2 |\phi_2\rangle\langle\phi_2| + m_3 |\phi_3\rangle\langle\phi_3|.
\]

The interaction operator corresponding to the bilinear term is
\[
\hat{I} = \sqrt{m_2 m_3}\left(e^{i\Delta\phi} |\phi_2\rangle\langle\phi_3|
+ e^{-i\Delta\phi} |\phi_3\rangle\langle\phi_2|\right).
\]
This is a Hermitian operator with matrix representation in the \(\{\phi_2,\phi_3\}\) basis:
\[
\hat{I} =
\sqrt{m_2 m_3}
\begin{pmatrix}
0 & e^{i\Delta\phi}\\
e^{-i\Delta\phi} & 0
\end{pmatrix}.
\]

\subsection{Eigenvalues and operator norm of \texorpdfstring{\(\hat{I}\)}{I}}

\begin{lemma}[Spectrum and norm of \(\hat{I}\)]
The operator \(\hat{I}\) has eigenvalues \(\lambda_{\pm} = \pm\sqrt{m_2 m_3}\), and its operator norm (spectral norm) satisfies
\[
\|\hat{I}\|_{\mathrm{op}} = \sqrt{m_2 m_3}.
\]
\end{lemma}

\begin{proof}
The characteristic polynomial of the \(2\times 2\) matrix
\[
\sqrt{m_2 m_3}
\begin{pmatrix}
0 & e^{i\Delta\phi}\\
e^{-i\Delta\phi} & 0
\end{pmatrix}
\]
is
\[
\det\left(
\sqrt{m_2 m_3}
\begin{pmatrix}
0 & e^{i\Delta\phi}\\
e^{-i\Delta\phi} & 0
\end{pmatrix} - \lambda I
\right)
= \det
\begin{pmatrix}
-\lambda & \sqrt{m_2m_3} e^{i\Delta\phi}\\
\sqrt{m_2m_3} e^{-i\Delta\phi} & -\lambda
\end{pmatrix}
= \lambda^2 - m_2 m_3.
\]
The eigenvalues are therefore \(\lambda_\pm = \pm\sqrt{m_2 m_3}\). As \(\hat{I}\) is Hermitian, its operator norm is the maximum absolute value of its eigenvalues, i.e.\ \(\|\hat{I}\|_{\mathrm{op}} = \sqrt{m_2 m_3}\).
\end{proof}

Note that the interaction scalar \(\sqrt{m_2 m_3}\cos\Delta\phi\) can be written as \(\frac{1}{2}\mathrm{Tr}(\hat{I} U)\) for an appropriate unitary \(U\) encoding the phase difference, but for our purposes, the norm bound is sufficient.

\subsection{Norm bounds for the multiplicity generator}

Consider the multiplicity generator \(\Sigma\) restricted to a bounded region of multiplicity space, e.g.\ \(m_2,m_3 \in [0,M]\) with \(M>0\).

\begin{proposition}[Bounds on \(\Sigma\)]
For \(m_2,m_3\in[0,M]\) and \(\Delta\phi\in\mathbb{R}\),
\[
\Sigma_{\min} \le \Sigma(m_2,m_3,\Delta\phi) \le \Sigma_{\max},
\]
where
\[
\Sigma_{\max} = M\log 2 + M\log 3 + |\beta| M,
\]
\[
\Sigma_{\min} = 0 \quad \text{if } \beta\ge 0;
\quad\text{or}\quad
\Sigma_{\min} = \min_{m_2,m_3\in[0,M]} \left(m_2\log 2 + m_3\log 3 - |\beta|\sqrt{m_2 m_3}\right)
\]
if \(\beta<0\). In particular, for \(\beta\ge 0\), \(\Sigma\) is bounded below by 0 and above by \(\Sigma_{\max}\) on the box \([0,M]^2\times\mathbb{R}\).
\end{proposition}

\begin{proof}
For \(\beta\ge 0\), we have
\[
m_2\log 2 + m_3\log 3 \ge 0,\quad
\sqrt{m_2 m_3}\cos\Delta\phi \ge -\sqrt{m_2 m_3},
\]
but since \(\beta\ge 0\) and \(\cos\Delta\phi\in[-1,1]\), the minimum of the interaction term is \(-\beta\sqrt{m_2m_3}\). However, at \(m_2=0\) or \(m_3=0\), this term vanishes, and the diagonal term is nonnegative. Thus the global minimum of \(\Sigma\) over \([0,M]^2\) is attained at \(m_2=m_3=0\), yielding \(\Sigma_{\min}=0\).

For the upper bound, we have
\[
m_2\log 2 + m_3\log 3 \le M\log 2 + M\log 3,
\]
and
\[
\beta\sqrt{m_2m_3}\cos\Delta\phi \le |\beta|\sqrt{m_2m_3} \le |\beta| M,
\]
since \(m_2,m_3\le M\). Hence \(\Sigma\le \Sigma_{\max}\) as claimed.

For \(\beta<0\), the interaction term can reduce \(\Sigma\) below zero, with the minimum attained at \(\cos\Delta\phi = 1\) and suitable \(m_2,m_3\); the stated expression gives the lower bound in that case.
\end{proof}

\subsection{Bounds on \texorpdfstring{\(\Phi\)}{Phi}}

Given the exponential form \(\Phi = \exp(\tilde\alpha\Sigma)\), we immediately obtain bounds on \(\Phi\) over bounded multiplicity domains. Assuming \(\tilde\alpha\ge 0\) and \(\beta\ge 0\), we have \(\Sigma\in[0,\Sigma_{\max}]\), hence
\[
1 \le \Phi \le \exp(\tilde\alpha \Sigma_{\max}).
\]

More generally, if \(\Sigma\in[\Sigma_{\min},\Sigma_{\max}]\), then
\[
\exp(\tilde\alpha\Sigma_{\min}) \le \Phi \le \exp(\tilde\alpha\Sigma_{\max}),\qquad
\tilde\alpha\in\mathbb{R}.
\]

These inequalities guarantee that, for bounded occupation numbers, the multiplicity renormalization factor \(\Phi\) remains finite and controlled, avoiding pathological blow-ups in physically relevant regimes.

\section{Small-Multiplicity Expansion}

For small multiplicities \(m_2,m_3\ll 1\), we can expand \(\Phi\) in powers of \(m_2,m_3\). Using \(\Sigma\) from \eqref{eq:sigma-appendix} and \(\Phi = e^{\tilde\alpha\Sigma}\), we obtain:
\[
\Phi = 1 + \tilde\alpha\Sigma + \frac{\tilde\alpha^2}{2}\Sigma^2 + O(\Sigma^3).
\]

To first order in \(m_2,m_3\), we can neglect the interaction term \(\beta\sqrt{m_2m_3}\) and quadratic contributions from \(\Sigma^2\); then
\[
\Phi \approx 1 + \tilde\alpha(m_2\log 2 + m_3\log 3).
\]

To second order, retaining the leading interaction term and the square of the diagonal term:
\[
\Phi \approx 1 + \tilde\alpha\left(m_2\log 2 + m_3\log 3 + \beta \sqrt{m_2 m_3}\cos\Delta\phi\right)
+ \frac{\tilde\alpha^2}{2} (m_2\log 2 + m_3\log 3)^2.
\]

This expansion shows explicitly how the prime-labeled diagonal multiplicity and coherent interaction appear as successive corrections to the geometric acceleration in the small-multiplicity regime.

\section{Summary of Mathematical Results}

We summarize the main formal statements established in this appendix:

\begin{itemize}
\item The prime-product logarithm \(\log\Pi(\mathbf{m}) = \sum_p m_p\log p\) is the unique linear functional (up to a constant factor) encoding multiplicative composition of prime-labeled multiplicity sectors.
\item The interaction term \(\sqrt{m_2 m_3}\cos\Delta\phi\) is the unique lowest-order real symmetric bilinear invariant satisfying natural symmetry, vanishing, and phase-dependence constraints.
\item The multiplicity renormalization factor \(\Phi\) must be exponential in an additive scalar invariant \(\Sigma\) if it factorizes multiplicatively under composition of independent multiplicity sectors.
\item The interaction operator \(\hat{I}\) associated with the coherent term has eigenvalues \(\pm\sqrt{m_2 m_3}\) and operator norm \(\sqrt{m_2 m_3}\).
\item On bounded multiplicity domains, the generator \(\Sigma\) and hence \(\Phi\) obey uniform bounds, ensuring controlled behaviour of the multiplicity-induced acceleration modifications.
\end{itemize}

These results provide a rigorous mathematical backbone for the two-prime multiplicity model and suggest straightforward generalizations to larger prime sets and more complex interaction structures.

\nocite{*}
\bibliographystyle{plain}
\bibliography{references}
\end{document}